\section{Results}
\label{sec:results}

Sections~\ref{sec:results-translation} and~\ref{sec:results-boundary} test
Proposition 1 and delimit its scope. Sections~\ref{sec:results-selection}
to~\ref{sec:results-aggregation} place this result within the pipeline
comparison in which it was found, and Section~\ref{sec:results-robustness}
summarises robustness controls that are reported in full in the supplementary
material. Unless stated otherwise, intervals are 95\% paired cluster bootstrap
intervals over target-test speakers and seeds.

\subsection{Exact validation equality in the matched translation pairs}
\label{sec:results-translation}

Table~\ref{tab:translation-audit} tests the stored consequences of
Proposition 1. Each cell of the audit is one matched pair of \texttt{none} and
\texttt{mean\_shift} runs that share the direction, seed, backbone, layer
setting, classifier, and every split, label-map, feature and search identifier.
The audit requires exactly one successful row per arm and rejects duplicate or
unmatched rows. No target score participates in the join or the inclusion
rule. For each classifier family and direction, the table counts the pairs
whose stored winning source-validation scores are exactly equal (Equal val.)
and whose selected classifier parameters are exactly equal (Equal params). It
also reports the mean target macro-F1 difference between \texttt{mean\_shift}
and \texttt{none} ($\Delta$), the number of pairs in which mean shift improves
the target score (Positive), and the chance baseline. The upper panel contains
the RBF SVM, for which Proposition 1 makes an exact prediction. The lower panel
contains the families whose fitting rules fall outside its assumptions
(Section~\ref{sec:results-boundary}). The configuration and input-ledger
hashes, all matched run identifiers and the unrounded values are released with
the audit.

\begin{table*}[!t]
\centering\small
\caption{Retrospective translation audit of matched \texttt{none}/\texttt{mean\_shift} pairs.}
\label{tab:translation-audit}
\begin{tabular}{llrrrrrr}
\toprule
Direction & Head & Pairs & Equal val. & Equal params & $\Delta$ & Positive & Chance \\
\midrule
\multicolumn{8}{l}{\textbf{Stationary kernel, Gram-determined fit: Proposition 1 applies}} \\
CREMA-D $\to$ RAVDESS & RBF SVM & 30 & 30 & 30 & +0.1969 & 30 & 0.1656 \\
RAVDESS $\to$ CREMA-D & RBF SVM & 30 & 30 & 30 & +0.1342 & 30 & 0.1665 \\
\midrule
\multicolumn{8}{l}{\textbf{Penalised intercept or iterative solver: outside those assumptions}} \\
CREMA-D $\to$ RAVDESS & Logistic & 30 & 0 & 22 & +0.1881 & 30 & 0.1656 \\
RAVDESS $\to$ CREMA-D & Logistic & 30 & 3 & 12 & +0.1442 & 30 & 0.1665 \\
CREMA-D $\to$ RAVDESS & Linear SVM & 30 & 0 & 19 & +0.1776 & 30 & 0.1656 \\
RAVDESS $\to$ CREMA-D & Linear SVM & 30 & 6 & 15 & +0.1456 & 30 & 0.1665 \\
CREMA-D $\to$ RAVDESS & MLP & 45 & 0 & 24 & +0.1701 & 45 & 0.1656 \\
RAVDESS $\to$ CREMA-D & MLP & 45 & 0 & 12 & +0.1287 & 45 & 0.1665 \\
\bottomrule
\end{tabular}
% Provenance: Filter: \texttt{grid-freeze-v3}, \texttt{blending=none}; matched \texttt{none}/\texttt{mean\_shift} rows only. Every family spans three backbones and five seeds; the MLP additionally has the learned \texttt{weighted} layer aggregation, hence its larger cell count. Generated by \texttt{tools/audit\_translation.py} from the configuration \texttt{configs/audit\_translation.yaml}; the JSON report retains input/configuration hashes and both run IDs for every pair.
\end{table*}


For the RBF SVM, every pair in both directions has exactly equal stored
validation scores and exactly equal selected hyperparameters. This is the
predicted equality, observed without exception for the classifier that the
proposition covers. Mean shift nevertheless improves target macro-F1 in every
pair. Source-side validation therefore assigned identical evidence to two
conditions whose target predictions differ, which is precisely the failure that
the proposition describes.

The ledger stores the winning validation score of each run, not the
predictions of every search candidate. The audit therefore verifies the
recorded winning scores and selected parameters rather than the entire search
surface, and the statement about every candidate remains analytic. Because the
cells share seeds and speakers, the table reports descriptive paired means
rather than a test that treats each cell as an independent replication. Mean
shift improves the target score in every matched pair of every family, but this
is a descriptive property of these cells against an unaligned baseline in one
corpus pair. It is not evidence that mean matching improves cross-corpus SER in
general, and no such claim is made.

\subsection{Classifiers outside the proposition's assumptions}
\label{sec:results-boundary}

Only the RBF SVM satisfies the fitting assumptions of Proposition 1. The lower
panel of Table~\ref{tab:translation-audit} reports the three families that the
assumptions exclude, as identified in advance in
Section~\ref{sec:translation-theory}. The linear SVM has a penalised synthetic
intercept, logistic regression uses a finite iterative solver, and the MLP
depends on initialisation, weight decay and early stopping.

These families behave as the assumptions predict. Exact validation equality
occurs in only a few forward pairs of logistic regression and the linear SVM,
in no reverse pair, and never for the MLP. The selected hyperparameters agree
in only part of the pairs. The MLP has more pairs per direction because it
additionally searches the learned \texttt{weighted} layer aggregation. None of
these results contradicts Proposition 1, whose statement does not extend to a
penalised intercept or to iterative neural training. These rows delimit the
proposition rather than test it, and they separate an exact identity between
objectives from the numerical behaviour of the solvers that approximate them.
The reduced-seed Transformer arm has no balanced matched cells and is excluded
by construction.

\subsection{Practical consequences of source-side selection}
\label{sec:results-selection}

The invariance in Section~\ref{sec:results-translation} is exact but local,
because it concerns two conditions that the criterion cannot separate.
Table~\ref{tab:headline} shows its consequence for the complete protocol by
comparing the configuration selected on source-side validation alone with the
best configuration present in the grid. As explained in
Section~\ref{sec:metrics}, the intervals in this table are $t$-intervals over
the five seeds rather than paired cluster bootstrap intervals.

\begin{table*}[!t]
  \centering
  \setlength{\tabcolsep}{5pt}
  \caption{Target macro-F1 of the source-validated and oracle configurations (mean over five seeds with 95\% $t$-intervals).}
  \label{tab:headline}
  \begin{tabular}{lrrrrr}
    \toprule
    pair & validated & oracle & gap & chance & majority \\
    \midrule
    RAVDESS $\rightarrow$ CREMA-D & 0.3156 [0.1880, 0.4432] & 0.4555 [0.4326, 0.4785] & 0.1399 [0.0231, 0.2568] & 0.1665 & 0.0425 \\
    CREMA-D $\rightarrow$ RAVDESS & 0.4994 [0.4592, 0.5395] & 0.5680 [0.5148, 0.6212] & 0.0687 [0.0090, 0.1283] & 0.1656 & 0.0444 \\
    \bottomrule
  \end{tabular}
  % Provenance: Filter: \texttt{freeze\_tag=grid-freeze-v3}, \texttt{blending=none}, 4986 runs. Mean over 5 seeds with a 95\% $t$-interval. Both baselines are analytic from the realised target-test priors.
\end{table*}


In both directions, the validated configuration exceeds the chance baseline on
average but falls short of the oracle, and the oracle gap is about twice as
large for RAVDESS $\rightarrow$ CREMA-D as for the reverse direction.
Figure~\ref{fig:selection} shows the gap for each seed. The forward direction
is where the protocol visibly fails. In this direction, \texttt{source\_val}
separates \texttt{none} ($0.7304$) from \texttt{zscore} ($0.7322$) by only
$0.0018$, whereas the same two rungs differ by $0.13$ on the target. Unlike the
source-translation result, this near-tie is empirical rather than an exact
invariance. In this case, the criterion is not symmetric but merely
uninformative at the available resolution. As a consequence, selection falls back to \texttt{none}
on two of the five forward seeds, and on one of them the selected configuration
scores \textbf{below the chance baseline}. The protocol uses no target labels
and gives no source-side indication of this outcome. In the reverse direction,
the protocol selects an aligned rung on all five seeds (\texttt{mean\_shift}
four times and \texttt{zscore} once), and its gap is correspondingly smaller.

\begin{figure*}[!t]
  \centering
  \includegraphics[width=\textwidth]{validated_vs_oracle.pdf}
  \caption{Validated and oracle target macro-F1 for each seed. Each validated
    bar is annotated with the rung selected by source-side validation, and the
    dotted line marks the chance baseline.}
  \label{fig:selection}
\end{figure*}

This result is limited to the selection of the alignment condition. Source-side
selection of the layer aggregation is assessed separately in
Section~\ref{sec:results-aggregation}.

\subsection{The affine comparison as pipeline context}
\label{sec:results-ladder}

The conditions among which the criterion selects form the affine ladder.
Table~\ref{tab:ladder} reports the source-selected target macro-F1 of each rung
together with candidate-row discrepancy summaries, and Table~\ref{tab:primary}
reports the pre-specified comparison family. The two parts of
Table~\ref{tab:ladder} use different averaging protocols. The discrepancy
columns average every inner setting, whereas the target scores use the
source-selected setting in each cell. Their juxtaposition is therefore
descriptive and is not a matched map-level correlation. The forward summaries
include the available Transformer cells, so they do not compare identical
candidate sets across seeds.

\begin{table*}[!t]
  \centering
  \caption{Source-selected target macro-F1 and mean marginal discrepancy of each alignment rung.}
  \label{tab:ladder}
  \begin{tabular}{lrrrrr}
    \toprule
    pair & rung & candidate rows & target macro-F1 & adaptive-own & reference-basis \\
    \midrule
    RAVDESS $\rightarrow$ CREMA-D & \texttt{none} & 153 & 0.2289 [0.2133, 0.2444] & 1432.93 & 36905.40 \\
     & \texttt{zscore} & 153 & 0.3599 [0.3483, 0.3726] & 51.79 & 480.85 \\
     & \texttt{mean\_shift} & 153 & 0.3595 [0.3442, 0.3762] & 104.05 & 19641.52 \\
     & \texttt{coral} & 828 & 0.3648 [0.3480, 0.3831] & 34.84 & 311.15 \\
     & \texttt{mkmmd\_diag} & 558 & 0.3497 [0.3347, 0.3663] & 67.34 & 1057.77 \\
     & \texttt{mkmmd\_full} & 558 & 0.3540 [0.3383, 0.3696] & 12.04 & 23.25 \\
    CREMA-D $\rightarrow$ RAVDESS & \texttt{none} & 135 & 0.2361 [0.1981, 0.2694] & 1045.24 & 23610.38 \\
     & \texttt{zscore} & 135 & 0.4321 [0.4070, 0.4547] & 41.00 & 331.64 \\
     & \texttt{mean\_shift} & 135 & 0.4178 [0.3931, 0.4386] & 76.60 & 8574.50 \\
     & \texttt{coral} & 810 & 0.4193 [0.3968, 0.4394] & 46.55 & 309.77 \\
     & \texttt{mkmmd\_diag} & 540 & 0.4212 [0.3984, 0.4417] & 46.37 & 830.54 \\
     & \texttt{mkmmd\_full} & 540 & 0.3885 [0.3653, 0.4099] & 10.71 & 17.71 \\
    \bottomrule
  \end{tabular}
  % Provenance: Filter: \texttt{freeze\_tag=grid-freeze-v3}, \texttt{blending=none}. Target intervals are a paired cluster bootstrap over target-test speakers and seeds, 2000 replicates; discrepancy columns are candidate-row means without intervals. Candidate counts differ because CORAL and MK-MMD have larger inner grids.
\end{table*}


Every aligned rung outperforms \texttt{none} in both directions. All ten of
these contrasts and both layer-aggregation contrasts belong to the
pre-specified primary family, and \textbf{all 14 survive Holm correction}, with
adjusted $p$-values at the bootstrap's resolution floor.

\begin{table*}[!t]
  \centering
  \caption{Pre-specified primary comparisons of target macro-F1 with Holm-corrected $p$-values.}
  \label{tab:primary}
  \begin{tabular}{lrrrrr}
    \toprule
    id & comparison & pair & difference in target macro-F1 & Holm $p$ & verdict \\
    \midrule
    L1 & \texttt{zscore - none} & RAVDESS $\rightarrow$ CREMA-D & +0.1310 [+0.1174, +0.1422] & $<$0.0070 & survives \\
    L1 & \texttt{zscore - none} & CREMA-D $\rightarrow$ RAVDESS & +0.1961 [+0.1501, +0.2449] & $<$0.0070 & survives \\
    L2 & \texttt{mean\_shift - none} & RAVDESS $\rightarrow$ CREMA-D & +0.1306 [+0.1159, +0.1439] & $<$0.0070 & survives \\
    L2 & \texttt{mean\_shift - none} & CREMA-D $\rightarrow$ RAVDESS & +0.1817 [+0.1354, +0.2302] & $<$0.0070 & survives \\
    L3 & \texttt{coral - none} & RAVDESS $\rightarrow$ CREMA-D & +0.1359 [+0.1253, +0.1476] & $<$0.0070 & survives \\
    L3 & \texttt{coral - none} & CREMA-D $\rightarrow$ RAVDESS & +0.1833 [+0.1373, +0.2293] & $<$0.0070 & survives \\
    L4 & \texttt{mkmmd\_diag - none} & RAVDESS $\rightarrow$ CREMA-D & +0.1208 [+0.1079, +0.1330] & $<$0.0070 & survives \\
    L4 & \texttt{mkmmd\_diag - none} & CREMA-D $\rightarrow$ RAVDESS & +0.1852 [+0.1406, +0.2317] & $<$0.0070 & survives \\
    L5 & \texttt{mkmmd\_full - none} & RAVDESS $\rightarrow$ CREMA-D & +0.1251 [+0.1149, +0.1351] & $<$0.0070 & survives \\
    L5 & \texttt{mkmmd\_full - none} & CREMA-D $\rightarrow$ RAVDESS & +0.1524 [+0.1057, +0.1991] & $<$0.0070 & survives \\
    A1 & \texttt{layer - last} & RAVDESS $\rightarrow$ CREMA-D & +0.0640 [+0.0541, +0.0748] & $<$0.0070 & survives \\
    A1 & \texttt{layer - last} & CREMA-D $\rightarrow$ RAVDESS & +0.1439 [+0.1140, +0.1746] & $<$0.0070 & survives \\
    A2 & \texttt{weighted - last} & RAVDESS $\rightarrow$ CREMA-D & +0.0597 [+0.0504, +0.0702] & $<$0.0070 & survives \\
    A2 & \texttt{weighted - last} & CREMA-D $\rightarrow$ RAVDESS & +0.1096 [+0.0888, +0.1326] & $<$0.0070 & survives \\
    \bottomrule
  \end{tabular}
  % Provenance: Paired cluster bootstrap over target-test speakers and seeds, 2000 replicates. $p$ values marked $<$ are at the bootstrap's resolution floor of $1/n_{\mathrm{boot}}$ and are not resolvable below it.
\end{table*}


The comparisons among aligned rungs are post hoc, and no evaluated condition is
shown to outperform \texttt{zscore}. To quantify the remaining uncertainty, we
bound the differences between each of the four other aligned conditions and
\texttt{zscore}, jointly across both directions. The eight-contrast Bonferroni
family uses the $99.375$th bootstrap percentile, targeting approximate 95\%
family-wise coverage. The largest corrected upper bound is $+0.0163$ forward
and $+0.0011$ reverse, both attained by CORAL, with point differences of
$+0.0049$ and $-0.0128$. Positive upper bounds are not evidence of an
advantage, and these results do not establish practical equivalence without a
predefined application-specific margin. The conditions can also differ in the
other direction. In the reverse direction, \texttt{mkmmd\_full} minus
\texttt{zscore} is $-0.0437$ $[-0.0559, -0.0320]$.

Two limits on reading these tables matter for the argument that follows.
First, the scores compare complete pipelines, including the classifier
geometries induced by the maps. By Lemma 2, the rungs differ in effective
regularisation and kernel geometry as well as in matched moments. The tables
therefore do not establish that covariance matching has no value or that
matching additional moments brings no benefit. Second, because of the
acceptance check described in Section~\ref{sec:ladder}, the MK-MMD rows
describe a finite procedure that includes its fallback, not a converged affine
MK-MMD solution.

\subsection{The discrepancy--transfer relationship is not measurement-invariant}
\label{sec:results-frames}

Alignment methods are commonly justified by the discrepancy they remove, which
presumes that the relationship between discrepancy and transfer is well
defined. Equation~\ref{eq:kernel-pullback} shows that this relationship is
defined only relative to a coordinate system. Table~\ref{tab:frames} and
Figure~\ref{fig:frames} measure the consequence on a 13-layer sweep of 2340
runs. For each direction, backbone, rung and seed, the Spearman correlation
between a layer's marginal discrepancy and its target macro-F1 is computed
across the 13 layers and then averaged over rungs and seeds. The pooled row
averages all 36 cells and five seeds. Because the cells share seeds and
features, these values are descriptive means rather than independent
replicates.

\begin{table*}[!t]
  \centering
  \caption{Spearman correlation between layer-wise marginal discrepancy and target macro-F1 under two measurement protocols.}
  \label{tab:frames}
  \begin{tabular}{lrrr}
    \toprule
    pair & backbone & $\rho$ (adaptive own) & $\rho$ (reference basis) \\
    \midrule
    RAVDESS $\rightarrow$ CREMA-D & \texttt{hubert} & -0.404 & 0.385 \\
     & \texttt{wav2vec2} & -0.372 & 0.224 \\
     & \texttt{wavlm} & -0.514 & 0.230 \\
    CREMA-D $\rightarrow$ RAVDESS & \texttt{hubert} & -0.450 & -0.038 \\
     & \texttt{wav2vec2} & -0.366 & 0.036 \\
     & \texttt{wavlm} & -0.560 & 0.012 \\
    \textbf{pooled} &  & \textbf{-0.444} & \textbf{0.142} \\
    \bottomrule
  \end{tabular}
  % Provenance: Filter: 13-layer sweep, 2340 runs, logreg, 6 rungs, 3 backbones, both directions, 5 seeds. Each correlation is computed across 13 layers within one direction, backbone, rung and seed, then averaged over rungs and seeds. The pooled row averages all 36 cells and 5 seeds.
\end{table*}


\begin{figure*}[!t]
  \centering
  \includegraphics[width=\textwidth]{frame_dependence.pdf}
  \caption{Measurement-protocol sensitivity. Left: Spearman $\rho$ between
    layer-wise marginal discrepancy and target macro-F1 for the 36
    (direction $\times$ backbone $\times$ rung) cells under the two protocols.
    Right: discrepancy along the CORAL shrinkage path, whose map limit is known
    analytically.}
  \label{fig:frames}
\end{figure*}

The pooled correlation is negative in each rung's adaptive own geometry and
positive in the reference-basis calculation. These opposite signs describe the
fixed sweep. Because the correlations share seeds and feature sets, an interval
that treats all cell-and-seed values as independent is not an inferential
justification for the sign change. Descriptively, 11 of the 36 individual cells
have uncorrected intervals on opposite sides of zero. This count is not used as
an independent significance claim.

The two statistics use the same features, layers and target scores and differ
only in their basis and bandwidth rule. \textbf{The correlation is therefore
not invariant to the stated measurement protocol.} The runs do not identify a
universally correct protocol. Instead, they show that a discrepancy--transfer
correlation must state its coordinate basis and bandwidth rule before it can be
compared with another report.

A second gap points in the same direction. The generalisation bound that
motivates treating discrepancy reduction as progress is stated over a
\emph{classifier-induced} divergence, defined on the symmetric difference of
the hypothesis class (Section~\ref{sec:related-alignment}). A kernel MMD at a
fixed bandwidth is not that quantity. It is a property of the two feature
samples alone and does not depend on the hypothesis class that the learner
searches. Reducing it is therefore not, in itself, a reduction of the term that
the theory bounds.

\paragraph{A map limit with known form.} One case fixes the interpretation
independently of any discrepancy measurement. With each covariance regularised
as $C + \varepsilon\,\mathrm{tr}(C)/d \cdot I$, the CORAL map satisfies
\begin{equation}
  M \;\longrightarrow\; \sqrt{\mathrm{tr}(C_t)/\mathrm{tr}(C_s)} \; I
  \label{eq:coral-limit}
\end{equation}
as $\varepsilon$ grows. The transform therefore collapses to a global scalar
rescaling plus a mean shift, that is, \texttt{mean\_shift} with one extra
degree of freedom. A shrinkage probe confirms this approach on real features.
The relative distance to the nearest scaled identity falls from $0.9254$ to
$0.0049$, and the fitted scalar converges to the predicted limit
(Supplementary Section~S5). The limiting map is a source-only transformation,
so the shared-scale invariance in Equation~\ref{eq:median-scale-invariance}
does not apply to it. Knowing the form of the map therefore does not fix the
measured discrepancy.

\paragraph{Recommendation.} We state the consequence as a recommendation
rather than as a characterisation of what other work reports, which we have not
surveyed. \textbf{A reported reduction in discrepancy should state the geometry
in which the discrepancy was measured, namely the kernel bandwidth and the
basis, and how each was fixed.} This matters when a data-dependent rule such as
the median heuristic sets the bandwidth, because the statistic then adapts to a
global rescaling of both samples (Equation~\ref{eq:median-scale-invariance}).
A source-defined fixed kernel is one explicit comparison protocol, not a
universal remedy, and its saturation under z-score is reported in
Supplementary Section~S4. An adaptive calculation answers a different question
and should be labelled with its basis and bandwidth rule.

\subsection{Cases where source validation remains informative}
\label{sec:results-aggregation}

Source-side validation is not uninformative in general, and the same grid
shows this. Replacing final-layer pooling with a fixed mid-stack layer (A1) or
with a learned softmax over all 13 states (A2) improves target macro-F1 in both
directions, and both comparisons survive correction (Table~\ref{tab:primary}).
These effects are larger than any difference \emph{among} the aligned rungs
($\le 0.0437$) and smaller than the effect of aligning at all. Usefulness on
one search axis does not imply usefulness on another. The same criterion that
separates layer settings cannot separate the two translation arms at all, so
the failure is specific to a decision rather than a property of source
validation as such.

\subsection{Robustness of the aligned-versus-unaligned contrast}
\label{sec:results-robustness}

The selection result rests on the comparison between aligned and unaligned
conditions, so we test whether this comparison depends on the label
harmonisation that the two corpora require. RAVDESS distinguishes
\texttt{calm} from \texttt{neutral}, whereas CREMA-D does not, and the main
analysis merges the two classes. Two pre-specified controls either drop RAVDESS
\texttt{calm} or additionally exclude \texttt{neutral} to give a five-class
task. Under both controls, every aligned-versus-\texttt{none} contrast remains
positive in both directions, with positive paired 95\% intervals and positive
one-sided lower bounds under a shared 12-contrast Bonferroni family
(Supplementary Section~S1). \textbf{The aligned-versus-unaligned conclusion is
therefore not an artefact of the calm-to-neutral mapping.}

Further analyses are reported in the supplement because they qualify the scope
of the results without supporting the central argument. Standard
label-shift corrections reduce macro-F1 in 238 of 240 (record, estimator)
pairs despite a small prior divergence. This bounds what a prior-only
correction can achieve in this setting, but it does not establish the
assumptions of these corrections for these corpora (Supplementary
Section~S2). Class-conditional discrepancy remains nonzero after alignment
under both label mappings. It is a descriptive diagnostic rather than an
estimate of $P(y \mid x)$, and it is not offered as an explanation of the
similarity among the aligned rungs (Supplementary Section~S4). Per-class F1 and
the confusion structure of the source-validated configuration are reported in
Supplementary Section~S6, and a historical blending probe that does not affect
the criterion audit is reported in Supplementary Section~S7.
